trying out 7 variations.
no proof: does not asks for mathetical consistency and proof
no constraints: removing all avoidance instruction of harmful optimization
minimal: just asking for optimized version in one sentence nothing else
no explanation: removes the need for explanation
no format: removes the instruction for structured output
no hardware: removes all the hardware specific guardrails
cot: adding a "think step by step about bottlenecks before writing code" in the instructions

We evaluate seven prompt variants, each removing or modifying a single component of the full prompt,
to isolate the contribution of individual design choices to optimization quality and reliability.

\paragraph{No Proof (\texttt{no\_proof}).}
The full prompt requires the model to provide a formal justification that the optimized kernel is
semantically equivalent to the original. Removing this requirement allows the model to apply
transformations without verifying correctness. This is consequential: on \texttt{shmembench},
the swap loop preserves an algebraic sum invariant, and without the proof requirement Gemini~2.5~Pro
exploits this to eliminate all shared memory operations entirely, producing a 213$\times$ apparent
speedup that does not reflect genuine hardware improvement. GPT-5 applies this transformation
regardless of variant.

\paragraph{No Constraints (\texttt{no\_constraints}).}
The full prompt includes explicit avoidance instructions, such as prohibitions against kernel fusion,
entry-point renaming, and modifications that sacrifice correctness for speed. Removing these
instructions gives the model unconstrained freedom to restructure the code. This can expose
creative optimizations that the constraints would otherwise suppress, but also risks producing
semantically incorrect or benchmark-invalidating optimizations.

\paragraph{Minimal (\texttt{minimal}).}
Rather than providing profiling context, roofline data, and structured instructions, this variant
reduces the prompt to a single sentence: \textit{``Optimize this CUDA kernel.''}
The model receives the kernel source with no information about where bottlenecks occur.
Without this context, the model must guess the optimization target, which frequently results
in degraded performance. On \texttt{md5hash}, GPT-5 under \texttt{minimal} produced code
51\% slower than the best-performing variant, consistent with optimizing the wrong bottleneck.

\paragraph{No Explanation (\texttt{no\_explanation}).}
The full prompt asks the model to accompany its code with a natural-language explanation of
each change and its rationale. Removing this requirement eliminates a form of implicit
chain-of-thought: articulating a justification encourages careful reasoning before committing
to a transformation. Without it, the model skips this reasoning step and proceeds directly
to code generation, which can increase the rate of subtle errors.

\paragraph{No Format (\texttt{no\_format}).}
The full prompt specifies a structured output format: optimized code in a fenced code block,
followed by a bulleted list of changes. Removing this specification causes the model to produce
free-form output, which may intermix explanation with code, embed code in prose, or omit
fenced delimiters entirely. This primarily affects the reliability of automated code extraction
rather than the quality of the optimization itself.

\paragraph{No Hardware (\texttt{no\_hardware}).}
The full prompt includes hardware-specific guardrails, such as the target compute capability,
shared memory limits per block, and a list of warp-level primitives unavailable on the target
architecture. Without these constraints, the model may generate code that is valid in principle
but incompatible with the deployment target. This is the primary cause of Gemini's compile
failures on \texttt{sobol} under \texttt{cot}: the model introduced \texttt{\_\_builtin\_ctz},
a host-only intrinsic, inside a \texttt{\_\_global\_\_} kernel, and could not correct the error
across three repair iterations.

\paragraph{Chain-of-Thought (\texttt{cot}).}
Unlike the other variants, \texttt{cot} adds rather than removes a prompt component.
A single instruction is prepended: \textit{``Think step by step about the performance bottlenecks
before writing any code.''}
This encourages the model to reason explicitly about profiling evidence prior to generating
the optimized kernel. Across both models and most benchmarks, \texttt{cot} is the most
consistently stable variant, producing results at or near the best observed performance with
no additional repair iterations relative to the full prompt.
